\section{Related Work}
\label{sec:related}

Prior attempts to build multi-dimensional game design frameworks fall into two categories:
structural taxonomies and empirical analyses. Elverdam and Aarseth~\cite{elverdam2007game}
identify 17 structural dimensions (space, time, player, control) but do not validate
them empirically or establish how they relate to player experience. Salen and
Zimmerman's~\cite{salen2003rules} foundational game design text articulates design
principles but does not propose a formal dimensional model. BGG's weight and rating system
provides the largest empirical dataset but reduces all design complexity to two numbers.

The TIGER framework is closest in spirit to the ``design lenses'' approach of Schell~\cite{schell2008art},
which proposes hundreds of named analytical perspectives on game design, but differs in
being empirically grounded: TIGER dimensions emerge from PCA on scored data rather than
from designer intuition. The closest computational precedent is Liapis et al.'s~\cite{liapis2013sentient}
design space mapping work in procedural level generation, which identifies dimensions
of level design space from generative exploration rather than from human scoring.
